\documentclass[../readings.tex]{subfiles}

\begin{document}
\subsection{Polynomial Interpolation}
\label{subsec:interpolation}

\textBox{%
Constructing a polynomial that passes exactly through a given set of data points.
}

\subsubsection{Existence and Uniqueness}

\addThm{Existence and Uniqueness}{%
Given $n+1$ distinct points $(x_i, y_i)$, there exists a unique polynomial $P_n(x)$ of degree at most $n$ such that $P_n(x_i) = y_i$ for all $i = 0, \dots, n$.
}

\subsubsection{Lagrange Form}

\textBox{%
Provides an explicit construction for the interpolating polynomial.
}

\addDef{Lagrange Basis Polynomials}{%
\begin{align*}
    L_i(x) = \prod_{j=0, j \neq i}^n \frac{x - x_j}{x_i - x_j}.
\end{align*}
The interpolating polynomial is $P_n(x) = \sum_{i=0}^n y_i L_i(x)$.
}

\addNote{%
\textbf{Computational Cost:}\enspace Evaluation of the Lagrange form costs $O(n^2)$ operations. Adding a new point requires recomputing all basis polynomials.
}

\subsubsection{Barycentric Interpolation}

\textBox{%
A numerically stable and efficient refinement of the Lagrange form.
}

\addDef{Barycentric Formula}{%
\begin{align*}
    P_n(x) = \frac{\sum_{i=0}^n \frac{w_i}{x - x_i} y_i}{\sum_{i=0}^n \frac{w_i}{x - x_i}}, \quad \text{where } w_i = \frac{1}{\prod_{j \neq i} (x_i - x_j)}.
\end{align*}
}

\addNote{%
\textbf{Advantages:}\enspace Evaluation costs only $O(n)$ once the weights $w_i$ are computed ($O(n^2)$). This is the standard recommended approach for general polynomial interpolation.
}

\subsubsection{Runge's Phenomenon}

\textBox{%
Equally spaced points are often a poor choice for high-degree interpolation.
}

\addNote{%
\textbf{The Runge Phenomenon:}\enspace For certain functions, such as $f(x) = 1/(1+25x^2)$ on $[-1, 1]$, the interpolation error near the boundaries increases exponentially as the degree $n \to \infty$. 
}

\addThm{Chebyshev Points}{%
The optimal nodes for minimizing interpolation error on $[-1, 1]$ are the roots of the Chebyshev polynomials:
\begin{align*}
    x_i = \cos\left( \frac{2i+1}{2n+2} \pi \right), \quad i = 0, \dots, n.
\end{align*}
Interpolation at Chebyshev points eliminates Runge's phenomenon and ensures near-optimal convergence.
}

\subsubsection{Cubic Splines}

\textBox{%
For large datasets, high-degree polynomial interpolation is avoided in favor of piecewise polynomials.
}

\addDef{Cubic Spline}{%
A piecewise cubic function $S(x)$ that is $C^2$ continuous (continuous values, first, and second derivatives) at the internal nodes.
}

\addNote{%
\textbf{Structure:}\enspace Computing the coefficients of a cubic spline requires solving a sparse tridiagonal linear system, which can be done in $O(n)$ time.
}

\subsubsection{Exercises}

\addExcr{%
\begin{enumerate}
    \item Construct the Lagrange form for points $(-1, 1), (0, 0), (1, 1)$. Verify it is $P_2(x) = x^2$.
    \item Use the barycentric formula to interpolate $f(x) = \sin(x)$ at 5 Chebyshev nodes.
    \item Compare the error plots for $f(x) = 1/(1+25x^2)$ using 15 equally spaced nodes vs. 15 Chebyshev nodes.
\end{enumerate}
}

\end{document}
